% AUTO-GENERATED by Github-Branch/tools/build_unified_source.ps1.
% Do not edit this file directly: update the transformation manifest instead.
\chapter{Numerical Methods and Monte Carlo Simulation}
\chapterstatus{Edited reuse from Team 3. Numerical outputs and source figures are excluded until parity and recomputation gates are closed.}
\sourceassetnote
\section{Definition and Fundamental Principles} 

Monte Carlo methods rely on the relationship between probability and measure. In mathematical terms, the probability of an event corresponds to the relative size of the subset of outcomes that represent it. The Monte Carlo approach reverses this idea: it estimates unknown quantities—such as volumes or integrals—by interpreting them as probabilities. Random samples are drawn from a known distribution, and the proportion of points falling within a region of interest provides a probabilistic estimate of its measure.

In general, the method allows the estimation of the expected value of a random variable through simulation. The principle is simple: use random sampling to approximate deterministic quantities.

Given a function \( f \) on the unit interval, the integral
\begin{equation} 
\alpha = \int_{0}^{1} f(x) \,dx
\end{equation}
can be written as an expected value \( \mathbb{E}[f(U)] \), where \( U \sim \mathcal{U}(0,1) \). If \( U_1, \ldots, U_n \) are independent uniform samples, the Monte Carlo estimator is
\begin{equation} 
\hat{\alpha}_n = \frac{1}{n} \sum_{i=1}^{n} f(U_i).
\end{equation}
By the Law of Large Numbers, \( \hat{\alpha}_n \to \alpha \) almost surely as \( n \to \infty \).

\begin{definition}[Monte Carlo Estimator]
Let \( X \) be a random variable with density \( \pi(x) \) and \( f \) a measurable function. The expected value
\[
\mathbb{E}_\pi[f(X)] = \int_{-\infty}^{+\infty} f(x)\pi(x)\,dx
\]
is approximated by
\[
\frac{1}{n} \sum_{i=1}^{n} f(X_i), \quad X_i \sim \pi(x), \; i.i.d.
\]
\end{definition}
\subsection{Mathematical Foundations} 

The theoretical justification of the Monte Carlo method is grounded in two classical probabilistic results: the Law of Large Numbers (LLN) and the Central Limit Theorem (CLT). The LLN ensures the consistency of the estimator, while the CLT quantifies its variability for finite sample sizes.

\begin{theorem}[Weak Law of Large Numbers]
Let \( \{X_i\}_{i \ge 1} \) be i.i.d. random variables with finite expectation \( \mu \). Then
\[
\bar{X}_n = \frac{1}{n} \sum_{i=1}^n X_i \xrightarrow{P} \mu,
\]
that is, the sample mean converges in probability to the expected value.
\end{theorem}

\begin{theorem}[Strong Law of Large Numbers]
If \( \mathbb{E}[|X_1|] < \infty \), then
\[
\frac{1}{n} \sum_{i=1}^n X_i \xrightarrow{\text{a.s.}} \mathbb{E}[X_1],
\]
meaning that the sample mean converges almost surely to the expected value.
\end{theorem}

\begin{remark}
The LLN guarantees that the Monte Carlo estimator is consistent under mild assumptions on \( f \). Increasing the number of samples \( n \) reduces the estimation error with probability approaching one, although the convergence remains asymptotic.
\end{remark}

\begin{theorem}[Central Limit Theorem]
Let \( X_1, X_2, \dots, X_n \) be i.i.d. random variables with mean \( \mu \) and finite variance \( \sigma^2 \). Then
\[
\frac{\bar{X}_n - \mu}{\sigma / \sqrt{n}} \xrightarrow{d} \mathcal{N}(0,1).
\]
Hence, for large \( n \), the sample mean is approximately normal with variance \( \sigma^2 / n \).
\end{theorem}

\begin{remark}
The CLT describes the distribution of the estimator for finite \( n \) and allows the construction of confidence intervals:
\[
\mathbb{P}\left( |\bar{X}_n - \mu| \leq \varepsilon \right) 
\approx \Phi\left( \frac{\sqrt{n}}{\sigma} \varepsilon \right),
\]
where \( \Phi(\cdot) \) denotes the standard normal cumulative distribution function.
\end{remark}
\subsection{Variance Estimation and Confidence Intervals }

Estimating the uncertainty associated with a Monte Carlo experiment is crucial for assessing the reliability of numerical results.  
Consider the integral of interest
\[
I = \mathbb{E}[f(X)] = \int f(x)\,p(x)\,dx,
\]
which we estimate by independent sampling $X_1,\dots,X_N$ from the distribution $p(x)$:
\[
\hat{I}_N = \frac{1}{N}\sum_{i=1}^N f(X_i).
\]
By the Central Limit Theorem, the sample mean $\hat{I}_N$ is asymptotically normal with variance $\mathrm{Var}[f(X)]/N$, that is
\[
\sqrt{N}\,(\hat{I}_N - I) \xrightarrow{d} \mathcal{N}(0,\,\sigma_f^2),
\]
where $\sigma_f^2 = \mathrm{Var}[f(X)]$. Since $\sigma_f^2$ is unknown, it is estimated through the sample variance:
\[
\hat{\sigma}_f^2 = \frac{1}{N-1}\sum_{i=1}^N \left(f(X_i) - \hat{I}_N\right)^2.
\]
This quantity represents the dispersion of the elementary results $f(X_i)$ and decreases with $1/N$, showing that the accuracy of Monte Carlo improves at the rate $O(N^{-1/2})$, regardless of the problem’s dimension.

The estimated variance allows for the construction of an asymptotic \emph{confidence interval} at level $1-\alpha$:
\[
\hat{I}_N \pm z_{\alpha/2}\,\frac{\hat{\sigma}_f}{\sqrt{N}},
\]
where $z_{\alpha/2}$ is the quantile of the standard normal distribution.  
This interval represents the region where the true value $I$ falls with approximate probability $1-\alpha$, providing an objective measure of the simulation’s precision.

A synthetic measure of the global error is the \emph{root mean square error} (RMSE), defined as
\[
\mathrm{RMSE}(\hat{I}_N) = \sqrt{\mathbb{E}\left[(\hat{I}_N - I)^2\right]} = \frac{\sigma_f}{\sqrt{N}}.
\]
This value measures the average quadratic deviation between the Monte Carlo estimate and the true expected value.  
The RMSE is the reference quantity for evaluating the quality of a method: halving the error requires four times as many simulations, highlighting the typical trade-off between computational cost and precision.

In summary, variance estimation, confidence intervals, and the RMSE constitute the statistical pillars that quantify the quality of Monte Carlo estimates, ensuring that empirical convergence is supported by a rigorous probabilistic foundation.
\subsection{Example: Monte Carlo Estimation of the Integral \texorpdfstring{$\int_0^1 e^{-x^2}\,dx$}{integral from 0 to 1 of exp(-x squared) dx}}

To illustrate the Monte Carlo estimation procedure and the construction of the confidence interval, consider the integral
\[
I = \int_0^1 e^{-x^2}\,dx,
\]
whose analytical value is known:
\[
I_{\text{analytic}} = \frac{\sqrt{\pi}}{2}\,\mathrm{erf}(1) \approx 0.746824.
\]
The goal is to estimate $I$ numerically via uniform sampling $x \sim U(0,1)$ and apply the sample mean formula:
\[
\hat{I}_N = \frac{1}{N}\sum_{i=1}^N e^{-X_i^2}.
\]
Each $X_i$ is drawn independently in $(0,1)$, and thus $\hat{I}_N$ is an unbiased estimator of $I$.  
By the Central Limit Theorem,
\[
\hat{I}_N \approx \mathcal{N}\!\left(I,\ \frac{\sigma_f^2}{N}\right),
\]
where the variance $\sigma_f^2 = \mathrm{Var}(e^{-X^2})$ is estimated empirically as
\[
\hat{\sigma}_f^2 = \frac{1}{N-1}\sum_{i=1}^N \left(e^{-X_i^2} - \hat{I}_N\right)^2.
\]
From this, the confidence interval at level $1-\alpha$ is obtained as
\[
\hat{I}_N \pm z_{\alpha/2}\frac{\hat{\sigma}_f}{\sqrt{N}}.
\]

\paragraph{Numerical Analysis.}
Assume we generate $N = 10^4$ samples (seed = 42). We obtain, for example,
\[
\hat{I}_{10^4} = 0.7510,\quad \hat{\sigma}_f = 0.1998,
\]
from which the 95\% confidence interval is
\[
0.7510 \pm 1.96\frac{0.1998}{\sqrt{10^4}} = [0.7467,\ 0.7554],
\]
which effectively contains the analytical value $I_{\text{analytic}}$.  
This demonstrates the statistical consistency of the Monte Carlo estimator and the validity of the asymptotic Gaussian approximation.

\paragraph{RMSE Behavior.}
In a repeated simulation with increasing numbers of samples $N$, the RMSE decreases according to the law
\[
\mathrm{RMSE}(\hat{I}_N) = \frac{\sigma_f}{\sqrt{N}},
\]
showing a linear trend in logarithmic scale:
\[
\log(\mathrm{RMSE}) \approx -\frac{1}{2}\log(N) + \text{constant}.
\]
This behavior highlights the slow yet steady convergence typical of the Monte Carlo method: doubling the precision requires quadrupling the number of simulations.

% Asset/code block omitted pending provenance acceptance.


\paragraph{Interpretation.}
The example of $\int_0^1 e^{-x^2}\,dx$ is paradigmatic: the integrand is smooth and free of discontinuities, so the variance is finite and the convergence perfectly follows the Central Limit Theorem.  
In financial contexts, the integral represents the expected value of a payoff $f(X)$ under a given distribution $p(x)$: the described procedure is therefore directly applicable to estimating the expected price of a derivative or a Monte Carlo risk measure.
\subsection{Application in Finance: European Option Pricing with Monte Carlo Simulation}

In this section, we present a first financial example illustrating how Monte Carlo simulation can be applied to option pricing.  
Specifically, we consider the valuation of a European call option under the classical Black--Scholes framework.  
The purpose here is to show how the stochastic structure of asset prices can be combined with statistical simulation to approximate expected payoffs.  
A more detailed discussion of the Black--Scholes model and its analytical properties is developed in a later section of this article (see Section~\sourceref{sec:XX} for further details).

Under the risk–neutral measure, the underlying asset price \(S(t)\) evolves according to the stochastic differential equation
\[
\frac{dS(t)}{S(t)} = r\,dt + \sigma\,dW(t),
\]
whose solution is
\[
S(T)=S(0)\exp\!\Big((r-\tfrac12\sigma^2)T+\sigma\sqrt{T}\,Z\Big), \qquad Z\sim\mathcal N(0,1).
\]
Here \(r\) is the continuously compounded risk–free interest rate, \(\sigma\) the volatility of the underlying asset, and \(W(t)\) a standard Brownian motion.

The fair (no–arbitrage) price of a European call option with strike \(K\) and maturity \(T\) is given by the discounted expectation of its terminal payoff:
\[
C = \mathbb E\!\left[e^{-rT}\,(S(T)-K)^+\right].
\]

Monte Carlo estimation of this expectation proceeds by simulating $N$ independent standard normal variables $Z_i$, mapping them into terminal prices $S_i(T)$, and computing the corresponding payoffs $C_i = e^{-rT}\max(S_i(T) - K, 0)$.  
The sample mean
\[
\widehat C_N = \frac{1}{N}\sum_{i=1}^N C_i
\]
is an unbiased and strongly consistent estimator of the true option price.  
Its standard error decreases at rate \(1/\sqrt{N}\), providing a quantitative measure of precision.  

For reference, the analytical Black--Scholes price of a call option is given by
\[
\text{BS}(S, \sigma, T, r, K) = S \, \Phi(d_1) - e^{-rT} K \, \Phi(d_2),
\]
where
\[
\begin{aligned}
d_1 &= \frac{\ln(S/K) + (r + \tfrac{1}{2}\sigma^2)T}{\sigma\sqrt{T}}, \qquad
d_2 = d_1 - \sigma\sqrt{T}.
\end{aligned}
\]

which serves as a benchmark against which the Monte Carlo estimate can be compared.

% Asset/code block omitted pending provenance acceptance.


\begin{remark}
This numerical experiment shows that the Monte Carlo estimate converges to the analytical Black--Scholes value as the number of simulations increases.  
The difference between the two prices can be interpreted as the sampling error, which can be reduced either by increasing $N$ or by applying variance–reduction techniques.  
A full theoretical discussion of the Black--Scholes model, its derivation, and extensions is provided in Section~\sourceref{sec:XX}.
\end{remark}

% ---- next approved source slice ----

\section{Random Number and Random Variable Generation} 

Monte Carlo simulation is one of the most powerful and versatile techniques for the numerical solution of complex problems in quantitative finance, engineering, and the applied sciences. 
At the core of any Monte Carlo method lies the generation of sequences of numbers that mimic independent samples drawn from one or more probability distributions. 
Although these sequences are produced deterministically by mathematical algorithms, they must possess statistical properties that make them \emph{indistinguishable}, under appropriate tests, from genuinely random sequences.

The first step in any simulation therefore consists in building a \emph{pseudo–random number generator} (RNG), i.e., an algorithm capable of producing a sequence of values that appear random and are uniformly distributed on the unit interval. 
From this uniform source, one can then obtain, via suitable transformations, samples from more complex distributions—such as the normal and the multivariate normal—widely used in pricing models and stochastic processes in finance.

The quality of Monte Carlo results depends directly on the quality of the generator used. 
A poorly designed sequence may introduce spurious correlations or fail to adequately cover the probability space, thereby compromising the statistical convergence of estimates. 
For this reason, the theory of pseudo–random generators blends modular arithmetic, number theory, and computational statistics, with the aim of producing sequences with long period, good multidimensional uniformity, and apparent independence.

In this section we present:
\begin{itemize}
    \item the fundamental principles of pseudo–random generators, with particular attention to the linear congruential generator (LCG);
    \item transformation methods that map uniform variates to other distributions of interest, such as the (uni)variate normal and the multivariate normal;
    \item the main tools for empirical validation of generator quality, including the $\chi^2$ test, QQ–plots, and the analysis of sample correlations;
    \item an application to pricing a European call using different RNGs, to compare their performance.
\end{itemize}

These concepts form the theoretical and practical foundation for reliable and reproducible Monte Carlo simulations, on which variance–reduction methods and quasi–Monte Carlo techniques—discussed in later sections—are built.
\subsection{Pseudo–Random Number Generators} 

The first step of any Monte Carlo simulation is the generation of sequences of numbers that mimic independent samples from a uniform distribution on the unit interval \([0,1]\). 
Since truly random numbers cannot be produced on a deterministic computer, we resort to \emph{pseudo–random generators}, i.e., algorithms that output deterministic sequences with statistical properties similar to those of independent random samples.

\begin{definition}[Linear Congruential Generator (LCG)]
A linear congruential generator produces a sequence of integers \( x_i \) and uniforms \( u_i = x_i / m \) via the recursion
\[
x_{i+1} = (a x_i + c) \bmod m, \qquad u_i = \frac{x_i}{m},
\]
where \( m \) is the \emph{modulus}, \( a \) the \emph{multiplier}, \( c \) the \emph{increment}, and \( x_0 \) the initial \emph{seed}. 
The sequence \(\{u_i\}\) is therefore deterministic but can, with appropriate parameters, display an empirical distribution close to uniform.
\end{definition}

The LCG is conceptually simple, computationally efficient, and easily reproducible. 
However, the quality of the sequence produced depends critically on the choice of the parameters \(a, c, m\). 
An inappropriate combination can produce very short cycles or regular patterns that compromise statistical performance.

% Asset/code block omitted pending provenance acceptance.


\begin{remark}[Period and reproducibility]
The generator’s \emph{period} is the length of the sequence before values repeat. 
If the parameters satisfy:
\[
\begin{gathered}
c \text{ and } m \text{ are coprime},\\
a - 1 \text{ is a multiple of all prime factors of } m,\\
a - 1 \text{ is a multiple of } 4 \text{ if } m \text{ is a multiple of } 4.
\end{gathered}
\]
the generator has \emph{full period}, equal to \(m\). 
This means it produces all possible values \(x_i \in \{0, 1, \ldots, m-1\}\) before repeating.
Moreover, fixing the same initial seed \(x_0\) yields exactly the same sequence, a key feature for verification and replication of experimental results.
\end{remark}

\begin{remark}[Empirical uniformity]
For an RNG to be usable in Monte Carlo, the sequence \(\{u_i\}\) should exhibit a uniform distribution on \([0,1]\). 
Formally, for any interval \(A \subset [0,1]\),
\[
\lim_{N \to \infty} \frac{1}{N}\sum_{i=1}^N \mathbf{1}_{\{u_i \in A\}} = \lambda(A),
\]
where \(\lambda(A)\) denotes Lebesgue measure (length) of \(A\).
In practice, this cannot be verified analytically and is assessed via statistical tests (e.g., the $\chi^2$ test) and graphical inspections of multidimensional uniformity.
\end{remark}

\begin{remark}[Apparent independence]
Another fundamental requirement is the absence of visible correlation between successive values. 
Since a deterministic generator cannot produce truly independent numbers, we require that any residual dependencies be \emph{statistically undetectable}. 
Equivalently, for any test function \(f\) and samples \(u_i, u_j\) sufficiently far apart,
\[
\mathrm{Cov}(f(u_i), f(u_j)) \approx 0.
\]
Good apparent independence ensures that simulations do not introduce systematic bias from the arithmetic structure of the generator.
\end{remark}

\begin{theorem}[Lattice structure of LCG sequences]
The consecutive \(d\)–tuples
\[
(u_i, u_{i+1}, \dots, u_{i+d-1})
\]
generated by an LCG lie on a finite number of parallel hyperplanes within the unit cube \([0,1]^d\). 
In particular, they can lie on at most \((d! \, m)^{1/d}\) distinct hyperplanes. 
This phenomenon, known as the \emph{lattice structure}, shows that LCG sequences—while appearing random in low dimension—exhibit geometric correlations in higher dimensions.
\end{theorem}

This limitation becomes significant in high–dimensional Monte Carlo problems, where projection onto regular planes reduces the effective coverage of the simulation space. 
To mitigate such effects, combined generators, such as those proposed by L’Ecuyer, combine multiple linear recurrences with different moduli. 
These generators retain the arithmetic simplicity of congruential models while achieving extremely long periods (up to about \(2^{319}\)) and improved multidimensional uniformity.

In summary, a good pseudo–random generator should:
\begin{itemize}
    \item possess a long period to avoid premature cycling;
    \item ensure empirically uniform distribution;
    \item exhibit no significant correlations among outputs;
    \item be efficient, portable, and reproducible.
\end{itemize}

The Linear Congruential Generator is thus the starting point for the modern theory of pseudo–random numbers, from which more sophisticated methods derive—such as combined generators or higher–order matrix recurrences (MRG32k3a). 
In the next sections, we show how such uniform sequences can be transformed into samples from arbitrary distributions, including the normal and multivariate normal, via specific transformation methods.
\subsection{Transformations from Uniforms to Arbitrary Distributions} 

Once a sequence of pseudo–random uniforms \( U_1, U_2, \dots \sim \mathcal{U}(0,1) \) is available,
one can generate samples from any other probability distribution
through suitable deterministic transformations. 
The key principle is that, if the cumulative distribution function (CDF) of the target distribution is known, one can invert the relationship between cumulative probabilities and variable values.

\begin{definition}[Inverse transform method]
Let \( F(x) \) be a continuous and strictly monotone CDF, and let \( U \sim \mathcal{U}(0,1) \).
Define
\[
X = F^{-1}(U),
\]
where \(F^{-1}\) is the quantile function associated with \(F\).
Then \(X\) has distribution \(F\).
\end{definition}

\noindent
The intuition is straightforward: a uniform \(U\) represents a random percentile between 0 and 1; 
applying the inverse CDF yields the corresponding quantile of the target distribution.

\begin{proof}[Proof]
For any \(x \in \mathbb{R}\):
\[
\mathbb{P}(X \leq x) = \mathbb{P}(F^{-1}(U) \leq x) 
= \mathbb{P}(U \leq F(x)) = F(x),
\]
since \(U\) is uniform on \((0,1)\). 
Hence \(X\) has CDF \(F\).
\end{proof}

\begin{remark}[Method properties]
The inverse transform is exact whenever \(F^{-1}\) is available in closed form. 
It requires only one uniform per sample and preserves monotonicity of the mapping \(U \mapsto X\), 
which is useful in \emph{variance reduction} techniques (e.g., the \emph{common random numbers} method).
\end{remark}

\begin{example}[Exponential distribution]
For an exponential distribution with mean \(\theta > 0\), 
the CDF is \( F(x) = 1 - e^{-x / \theta} \) for \(x \geq 0\). 
Inverting yields:
\[
F^{-1}(u) = -\theta \ln(1 - u),
\]
and for \(U \sim \mathcal{U}(0,1)\),
\[
X = -\theta \ln U.
\]
This illustrates the simplicity and efficiency of the method when a closed–form inverse exists.
\end{example}

% Asset/code block omitted pending provenance acceptance.


\bigskip
\noindent
In many financial and physical applications, however, the standard normal distribution \(\mathcal{N}(0,1)\) is central. 
Since the normal CDF \(\Phi(x)\) has no closed–form inverse, alternative transformations are used, the most famous being the Box–Muller transform.

\begin{definition}[Box–Muller transform]
Let \( U_1, U_2 \) be independent uniforms on \((0,1)\). 
Define:
\[
Z_1 = \sqrt{-2 \ln U_1}\cos(2\pi U_2), \qquad
Z_2 = \sqrt{-2 \ln U_1}\sin(2\pi U_2).
\]
Then \( Z_1, Z_2 \sim \mathcal{N}(0,1) \) are independent standard normals.
\end{definition}

\begin{proof}[Proof idea]
The method relies on the polar representation of the bivariate normal density:
\[
f(z_1, z_2) = \frac{1}{2\pi} \exp\!\left(-\frac{z_1^2 + z_2^2}{2}\right).
\]
The radius \( R = \sqrt{z_1^2 + z_2^2} \) has density \( f_R(r) = r e^{-r^2/2} \) for \( r > 0 \),
while the angle \(\Theta\) is uniform on \([0, 2\pi)\). 
Applying the inverse transform to \(R\) yields
\[
R = \sqrt{-2 \ln U_1}, \qquad \Theta = 2\pi U_2,
\]
and inverting the coordinates gives \( (Z_1, Z_2) = (R\cos\Theta, R\sin\Theta) \).
\end{proof}

% Asset/code block omitted pending provenance acceptance.



\begin{remark}[Practical considerations]
Box–Muller is exact and generates two normals per uniform pair. 
However, trigonometric evaluations can be costly at scale. 
Hence, faster variants like the Marsaglia–Bray algorithm generate points uniformly in the unit disk and rescale them to obtain \( (Z_1, Z_2) \),
avoiding sine and cosine.
\end{remark}

\begin{remark}[Inverse transform for the normal]
An alternative is to apply the inverse transform using the standard normal CDF:
\[
Z = \Phi^{-1}(U).
\]
Since \(\Phi^{-1}\) has no closed form, highly accurate numerical approximations (e.g., Beasley–Springer–Moro or Acklam) are used,
achieving errors below \(10^{-9}\).
This approach is often preferred in computational finance, where deterministic input–output mapping 
aids reproducibility and sensitivity analysis.
\end{remark}

\bigskip
In summary, the inverse transform and Box–Muller provide the two fundamental paradigms 
for generating non–uniform variates. 
The former is conceptually universal but requires access to \(F^{-1}\);
the latter is specific to the normal distribution yet extremely precise and reproducible. 
Both underlie more complex sampling methods, such as multivariate distributions or correlated transformations, 
addressed next.
\subsection{Generating Multivariate Normal Variates} 

Many financial applications require the simultaneous simulation of multiple correlated random variables, such as the returns of different assets in a portfolio or the components of a vector of risk factors. 
In such cases it is common to assume the random vector follows a multivariate normal distribution, which generalizes the univariate normal by incorporating cross–correlations.

\begin{definition}[Multivariate normal distribution]
A random vector \( X \in \mathbb{R}^d \) has a multivariate normal distribution with mean \(\mu \in \mathbb{R}^d\) and covariance matrix \(\Sigma \in \mathbb{R}^{d \times d}\), denoted
\[
X \sim \mathcal{N}(\mu, \Sigma),
\]
if for every \(a \in \mathbb{R}^d\) the linear combination \(a^\top X\) is univariate normal with
\[
a^\top X \sim \mathcal{N}(a^\top \mu,\, a^\top \Sigma a).
\]
Its density is
\[
f_X(x) = \frac{1}{(2\pi)^{d/2}|\Sigma|^{1/2}}
\exp\!\left(-\frac{1}{2}(x - \mu)^\top \Sigma^{-1} (x - \mu)\right).
\]
\end{definition}

\noindent
Sampling \( X \sim \mathcal{N}(\mu, \Sigma) \) can be reduced to the standard case \( Z \sim \mathcal{N}(0, I_d) \) via a linear transformation that introduces the desired covariance.

\begin{theorem}[Sampling via Cholesky factorization]
Let \(\Sigma\) be symmetric positive definite.
Then there exists a unique lower–triangular matrix \(L\) such that
\[
\Sigma = L L^\top.
\]
If \(Z \sim \mathcal{N}(0, I_d)\) is a vector of independent standard normals, then
\[
X = \mu + LZ
\]
has distribution \(\mathcal{N}(\mu, \Sigma)\).
\end{theorem}

\begin{proof}
We have
\[
\mathbb{E}[X] = \mu + L\,\mathbb{E}[Z] = \mu, \qquad
\mathrm{Cov}[X] = L\,\mathrm{Cov}[Z]\,L^\top = L I_d L^\top = LL^\top = \Sigma.
\]
A linear image of a normal vector is normal, proving the claim.
\end{proof}

% Asset/code block omitted pending provenance acceptance.


\begin{remark}[Geometric interpretation]
Cholesky acts as a rotation–scaling of the standard normal space \(Z\) 
to obtain the dependence structure encoded by \(\Sigma\).
Each row of \(L\) specifies a linear combination of the components of \(Z\), 
thus introducing correlations among the simulated variables.
\end{remark}

\begin{remark}[Computational efficiency]
Cholesky is numerically efficient, requiring about \(O(d^3 / 3)\) operations, and can be precomputed once if \(\Sigma\) is fixed. 
It is stable for well–conditioned matrices; 
when correlations are near perfect (i.e., \(\Sigma\) is nearly singular), 
alternatives such as eigen–decomposition or \emph{singular value decomposition (SVD)} may be preferable.
\end{remark}

\begin{remark}[Alternatives to Cholesky]
An equivalent route is the spectral decomposition
\[
\Sigma = V \Lambda V^\top,
\]
where \(V\) contains eigenvectors and \(\Lambda\) is the diagonal matrix of eigenvalues. 
Setting \(A = V \Lambda^{1/2}\) yields \(X = \mu + A Z\) with covariance \(\Sigma\). 
This is useful if \(\Sigma\) is not strictly positive definite or when explicit control of principal directions is desired.
\end{remark}

\bigskip
\noindent
Implementation–wise, the simulation proceeds as follows:
\begin{enumerate}
    \item generate \(Z = (Z_1, \dots, Z_d)\) with independent \(Z_i \sim \mathcal{N}(0,1)\);
    \item compute the Cholesky factorization \(\Sigma = LL^\top\);
    \item obtain the correlated vector \(X = \mu + LZ\).
\end{enumerate}

This approach is simple and general: 
once a good sample of independent normal variates is available, 
correlations can be imposed directly and controllably via the covariance matrix. 
It is widely used in computational finance, e.g.:
\begin{itemize}
    \item simulating multi–asset portfolios in multivariate Black–Scholes models;
    \item generating risk scenarios for Value–at–Risk or Expected Shortfall;
    \item constructing trajectories for correlated stochastic processes, such as multidimensional Brownian motions.
\end{itemize}

\begin{example}[Simulating two correlated assets]
Let \(S_1, S_2\) be assets whose log–returns follow a bivariate normal with zero mean, 
standard deviations \(\sigma_1, \sigma_2\), and correlation \(\rho\).
The covariance matrix is
\[
\Sigma = 
\begin{pmatrix}
\sigma_1^2 & \rho\,\sigma_1\sigma_2 \\
\rho\,\sigma_1\sigma_2 & \sigma_2^2
\end{pmatrix}.
\]
A Cholesky factor \(L\) is:
\[
L =
\begin{pmatrix}
\sigma_1 & 0 \\
\rho\,\sigma_2 & \sigma_2\sqrt{1-\rho^2}
\end{pmatrix}.
\]
Given \(Z = (Z_1, Z_2)^\top \sim \mathcal{N}(0, I_2)\), 
we obtain
\[
\begin{pmatrix} X_1 \\ X_2 \end{pmatrix}
= LZ =
\begin{pmatrix}
\sigma_1 Z_1 \\
\rho\,\sigma_2 Z_1 + \sigma_2 \sqrt{1-\rho^2}\,Z_2
\end{pmatrix},
\]
which ensures \(\mathrm{Corr}(X_1, X_2) = \rho\), as desired.
\end{example}

\bigskip
Generating multivariate normals is thus a key step 
in building realistic Monte Carlo models. 
It allows the dependence structure observed in historical data to be incorporated 
and underpins advanced techniques such as simulation of correlated stochastic processes 
and generation of financial scenarios consistent with multivariate risk models.
\paragraph{Principal Component Analysis (PCA)}
\begin{remark}[PCA as an alternative to Cholesky]
An alternative to Cholesky for generating multivariate normals is \emph{Principal Component Analysis (PCA)}. 
Again we start from the spectral decomposition of the covariance matrix:
\[
\Sigma = V \Lambda V^\top,
\]
where \(V = [v_1, \dots, v_d]\) is the orthogonal matrix of eigenvectors and \(\Lambda = \mathrm{diag}(\lambda_1, \dots, \lambda_d)\) is the diagonal matrix of nonnegative eigenvalues, conventionally ordered \(\lambda_1 \geq \lambda_2 \geq \dots \geq \lambda_d\).

The transformation
\[
X = \mu + V \Lambda^{1/2} Z, \qquad Z \sim \mathcal{N}(0, I_d),
\]
produces \(X \sim \mathcal{N}(\mu, \Sigma)\), as in the Cholesky case.
PCA offers an interpretive advantage: 
the directions \(v_i\) are the \emph{principal components}, i.e., orthogonal directions along which variance is maximized, with eigenvalues \(\lambda_i\) quantifying the variance carried by each component.
\end{remark}

\begin{remark}[Dimensionality reduction]
In many high–dimensional simulations, 
most of the total variance of \(\Sigma\) is concentrated in the first few eigenvalues.
Then one may approximate \(X\) by retaining only the first \(k < d\) principal components:
\[
X^{(k)} = \mu + \sum_{i=1}^{k} \sqrt{\lambda_i}\, v_i Z_i,
\]
with independent \(Z_i \sim \mathcal{N}(0,1)\).
This approximation minimizes the mean squared error
\[
\mathbb{E}\!\left[\|X - X^{(k)}\|^2\right] = \sum_{i=k+1}^{d} \lambda_i,
\]
i.e., the “discarded” variance from omitted components.
\end{remark}

\begin{remark}[Applications in quantitative finance]
PCA is widely used in high–dimensional financial simulation, 
as it reduces the number of risk drivers by simulating only principal components.
Applications include:
\begin{itemize}
    \item simulating term structures (e.g., Heath–Jarrow–Morton, Libor Market Model), 
          where a few factors explain most of the volatility of forward rates;
    \item dimensionality reduction in multi–asset portfolio simulation, 
          retaining only the leading systematic factors;
    \item compressing correlation or covariance data to accelerate pricing or stress–testing.
\end{itemize}
PCA also gives a clear geometric view of correlation structure: 
principal eigenvectors identify dominant risk directions, 
while eigenvalues quantify the variance carried by each factor.
\end{remark}

\begin{example}[PCA approximation]
Consider an empirical covariance matrix \(\Sigma \in \mathbb{R}^{5 \times 5}\) estimated on five asset returns.
Suppose the first two eigenvalues account for 85\% of total variance.
Then, generating only \(Z_1, Z_2 \sim \mathcal{N}(0,1)\) and setting
\[
X^{(2)} = \mu + \sqrt{\lambda_1}\, v_1 Z_1 + \sqrt{\lambda_2}\, v_2 Z_2,
\]
yields a two–dimensional simulation that captures most of the portfolio’s risk structure, 
significantly reducing computational cost.
\end{example}

\bigskip
\noindent
In conclusion, PCA is an orthogonal counterpart to Cholesky:
while the latter preserves the algebraic structure of \(\Sigma\) and is numerically efficient, 
PCA makes the geometry explicit and enables parsimonious representations of dominant risk factors. 
Both are fundamental for efficient simulation of correlated variates in multivariate Monte Carlo models.
\subsection{Empirical Validation of Generators} 

The statistical quality of a pseudo–random generator cannot be judged solely from its theoretical formulation: 
its uniformity, independence, and correlation properties must be empirically verified. 
The main methodologies include the \emph{uniformity $\chi^2$ test}, the \emph{QQ–plot} analysis, and the \emph{verification of empirical correlations}.

\begin{remark}[Purpose of empirical checks]
An ideal generator should produce a sequence \(\{U_i\}\) that:
\begin{enumerate}
    \item is uniformly distributed over \([0,1]\);
    \item exhibits apparent independence between successive terms;
    \item shows no spurious correlations or geometric patterns in multidimensional samples.
\end{enumerate}
Since a deterministic generator cannot satisfy these conditions in the absolute sense, deviations from the theoretical assumptions are required to be statistically negligible within test tolerances.
\end{remark}
\paragraph{Uniformity $\chi^2$ test}

The $\chi^2$ test is the classical method to check whether observations \(U_1, \dots, U_N\) follow a uniform distribution on \([0,1]\).
Partition \([0,1]\) into \(k\) equal subintervals and compare the observed frequencies \(O_j\) with the expected frequencies \(E_j = N/k\).

\begin{definition}[Test statistic]
The test statistic is
\[
\chi^2 = \sum_{j=1}^{k} \frac{(O_j - E_j)^2}{E_j}.
\]
Under the null hypothesis of uniformity and for large \(N\), the statistic is approximately \(\chi^2_{k-1}\)–distributed with \(k-1\) degrees of freedom.
\end{definition}

\noindent
A significantly large statistic indicates that observed frequencies differ from theoretical ones beyond what random fluctuations would suggest. 
In practice, the test is applied for several values of \(k\) and on partial sequences to assess local uniformity as well.

\begin{remark}[Interpretation]
If the p–value is too small (e.g., \(p < 0.01\)), 
the uniformity hypothesis is rejected: 
this may indicate that the generator favors certain parts of the interval or that the modulus arithmetic induces periodic patterns. 
An excessively high p–value (\(p \approx 1\)) may also suggest over–regularity, a sign of imperfect pseudo–randomness.
\end{remark}

% Asset/code block omitted pending provenance acceptance.
\paragraph{QQ–plot analysis for normality and other distributions}

When generating transformed variables (e.g., via Box–Muller or inverse CDF), 
it is appropriate to check whether the empirical distribution matches the target one.
The \emph{QQ–plot (Quantile–Quantile plot)} provides a graphical comparison.

\begin{definition}[Constructing the QQ–plot]
Given samples \(X_1, \dots, X_N\) and a theoretical distribution \(F\),
sort the data \(X_{(1)} \le \dots \le X_{(N)}\) and take the theoretical quantiles
\[
q_i = F^{-1}\!\left(\frac{i - 0.5}{N}\right), \qquad i = 1, \dots, N.
\]
The scatter of points \((q_i, X_{(i)})\) should align along the bisector \(y = x\) if the empirical and theoretical distributions agree.
\end{definition}

\begin{remark}[Interpreting QQ–plots]
Systematic deviations from the bisector signal differences between empirical and theoretical distributions:
\begin{itemize}
    \item curvature indicates heavier or lighter tails than the theoretical reference;
    \item slopes different from 1 indicate variance mismatch;
    \item offsets in central quantiles indicate mean bias.
\end{itemize}
QQ–plots are simple yet powerful diagnostics for validating normal or uniform generators and for comparing different transformation methods.
\end{remark}

% Asset/code block omitted pending provenance acceptance.
\paragraph{Checking empirical correlations}

For multidimensional generators or time sequences, further quality assessment comes from empirical correlations.

\begin{definition}[Empirical correlation]
Given samples \(X^{(1)}, \dots, X^{(N)} \in \mathbb{R}^d\), the empirical correlation between components \(i\) and \(j\) is
\[
\rho_{ij}^{\mathrm{emp}} = 
\frac{\mathrm{Cov}_{\mathrm{emp}}(X_i, X_j)}{\sqrt{\mathrm{Var}_{\mathrm{emp}}(X_i)\,\mathrm{Var}_{\mathrm{emp}}(X_j)}}.
\]
If samples are generated with theoretical covariance \(\Sigma\), 
we expect \(\rho_{ij}^{\mathrm{emp}} \approx \Sigma_{ij}/\sqrt{\Sigma_{ii}\Sigma_{jj}}\)
within statistical fluctuations of order \(O(N^{-1/2})\).
\end{definition}

\noindent
This check is crucial in multivariate models, 
where dependence must be respected not only on average but also sample–wise.
Correlation analysis can also reveal anomalies due to poor numerical conditioning in Cholesky or rounding issues.

\begin{remark}[Sampling uncertainty]
Because empirical correlations are noisy, 
it is useful to compare the sample value with a confidence interval approximated by
\[
\rho_{ij}^{\mathrm{emp}} \pm z_{\alpha/2}\sqrt{\frac{1 - (\rho_{ij}^{\mathrm{emp}})^2}{N-2}},
\]
where \(z_{\alpha/2}\) is the standard normal quantile for the desired confidence level.
\end{remark}

% Asset/code block omitted pending provenance acceptance.
\paragraph{Application: comparing RNGs via European call pricing}

A practical way to compare generators is to compute, with each RNG, the Monte Carlo price of a European call:
\[
C_0 = e^{-rT}\,\mathbb{E}\!\left[(S_T - K)^+\right], \qquad 
S_T = S_0 e^{(r - \frac{1}{2}\sigma^2)T + \sigma \sqrt{T}\,Z}.
\]
Comparing estimates from different RNGs (e.g., LCG, Mersenne Twister, Sobol, Halton) in terms of standard error and sample variance
helps assess numerical stability and statistical efficiency.
A high–quality RNG should produce consistent results within the estimated uncertainty, independent of the specific realization.

\bigskip
\noindent
\textbf{Summary.}  
The $\chi^2$ test, QQ–plots, and empirical correlation checks are fundamental tools for validating pseudo–random generators. 
They ensure that sequences used in Monte Carlo simulations are statistically reliable, 
so that discrepancies from theory arise only from sampling noise rather than intrinsic flaws in the generator.

% ---- next approved source slice ----

\section{Path Simulation and Discretization: Brownian, GBM, and Jump–Diffusion} 


Building on the random number generation and transformation techniques introduced in the previous section,  
this section presents practical methods for simulating stochastic paths that govern asset price dynamics.  
We focus on one– and multi–dimensional Brownian motion (BM), geometric Brownian motion (GBM), and their extensions to jump–diffusion (JD) processes, which introduce discontinuities and heavy tails.  
These models form the numerical core of Monte Carlo valuation methods in quantitative finance, providing the basis for path–dependent pricing, risk analysis, and scenario generation.  
Formal definitions and theoretical properties are discussed in Section~\sourceref{sec:??}.
\subsection{One-Dimensional Brownian Motion}

A \textbf{standard Brownian motion} $W(t)$ on $[0,T]$ is a stochastic process satisfying the following fundamental properties:
\begin{enumerate}
    \item $W(0) = 0$;
    \item The trajectories are continuous with probability $1$;
    \item The increments are independent;
    \item The increments are Gaussian: $W(t) - W(s) \sim \mathcal{N}(0, t-s)$ for $0 \le s < t$.
\end{enumerate}

Hence, $W(t) \sim \mathcal{N}(0, t)$.  
A \textbf{Brownian motion with drift and constant volatility} is defined as
\[
X(t) = \mu t + \sigma W(t),
\]
which satisfies the stochastic differential equation
\[
dX(t) = \mu\,dt + \sigma\,dW(t),
\]
and has the marginal distribution $X(t) \sim \mathcal{N}(\mu t, \sigma^2 t)$.

In the case of time-dependent coefficients $\mu(t)$ and $\sigma(t)$, the general solution is given by
\[
X(t) = X(0) + \int_0^t \mu(s)\,ds + \int_0^t \sigma(s)\,dW(s).
\]
\paragraph{Random Walk Construction}

The most direct construction of a discrete Brownian motion consists in simulating the points
\[
(W(t_1), W(t_2), \dots, W(t_n)), \quad 0 < t_1 < t_2 < \dots < t_n = T,
\]
using the property of independent increments:
\[
W(t_{i+1}) = W(t_i) + \sqrt{t_{i+1}-t_i}\,Z_{i+1}, \qquad Z_i \sim \mathcal{N}(0,1).
\]

For a process with drift $\mu$ and volatility $\sigma$, we have:
\[
X(t_{i+1}) = X(t_i) + \mu (t_{i+1}-t_i) + \sigma \sqrt{t_{i+1}-t_i}\,Z_{i+1}.
\]

This construction is \emph{exact} at the discrete grid points and forms the basis for the Euler method used in stochastic differential equations (see Section~\sourceref{sec:disc}).
\paragraph{Brownian Bridge Construction}

An alternative to the sequential simulation is the \textbf{Brownian bridge construction}, which generates the process values by conditioning on the known endpoints.

Let $W(u)=x$ and $W(t)=y$. Then, for $s \in (u,t)$:
\[
W(s) \mid (W(u)=x, W(t)=y) \sim 
\mathcal{N}\!\left(
\frac{(t-s)x + (s-u)y}{t-u},\;
\frac{(s-u)(t-s)}{t-u}
\right).
\]

The intermediate value is therefore the sum of a deterministic component (a linear interpolation) and a stochastic component (Gaussian noise).  
This method produces trajectories with the same marginal distribution as the random walk construction, but the generation order — from the endpoints toward the midpoints — allows for a reduction in variance when using low-discrepancy sequences or \emph{variance reduction} techniques.
\paragraph{PCA (Principal Component Analysis) Construction}

Let $W = (W(t_1), W(t_2), \dots, W(t_n))^\top$ denote the discretized Brownian motion vector. Its covariance matrix is
\[
C_{ij} = \mathbb{E}[W(t_i) W(t_j)] = \min(t_i, t_j).
\]
We seek a matrix $A$ such that $A A^\top = C$.  
The eigenvalue decomposition of $C$ yields
\[
C v_k = \lambda_k v_k, \qquad k = 1, \dots, n,
\]
which leads to
\[
W = A Z = \sum_{k=1}^n \sqrt{\lambda_k}\, v_k Z_k, \qquad Z_k \sim \mathcal{N}(0,1).
\]

The first few principal components (associated with the largest eigenvalues) explain most of the path variance.  
\paragraph{Multidimensional Extension}

For a $d$-dimensional Brownian motion $X(t) \in \mathbb{R}^d$:
\[
X(t) = \mu t + B W(t),
\]
where $W(t)$ is a standard Brownian motion in $\mathbb{R}^d$ and $B$ satisfies $B B^\top = \Sigma$.  
The increments satisfy
\[
X(t) - X(s) \sim \mathcal{N}\!\big((t-s)\mu, (t-s)\Sigma\big).
\]

All the previous constructions extend naturally to this multidimensional setting:
\begin{itemize}
    \item The \emph{random walk} construction applies independently to each coordinate;
    \item The \emph{Brownian bridge} can be applied by conditioning simultaneously on all coordinates;
    \item The \emph{PCA construction} uses the full covariance decomposition $(C \otimes \Sigma)$, where $C_{ij} = \min(t_i,t_j)$.
\end{itemize}

\subsection{Geometric Brownian Motion and Asset Simulation} 

Geometric Brownian motion (GBM) is the standard model for the stochastic evolution of asset prices in continuous time.  
It extends Brownian motion by ensuring positivity of prices and modeling proportional---rather than absolute---changes in value.  
The process provides the foundation of the Black--Scholes framework and of most Monte Carlo pricing engines.
\paragraph{Model definition}

A process $S(t)$ is called a \textbf{geometric Brownian motion} (GBM) if it satisfies the following stochastic differential equation:
\[
dS(t) = \mu S(t)\,dt + \sigma S(t)\,dW(t),
\]
where $\mu$ is the mean growth rate (\emph{drift}), $\sigma$ the volatility, and $W(t)$ a standard Brownian motion.

Applying Itô’s lemma to the logarithm of the price yields:
\[
d(\log S(t)) = \left(\mu - \frac{1}{2}\sigma^2\right)dt + \sigma\,dW(t),
\]
and integrating with respect to time gives:
\[
S(t) = S(0)\exp\!\left[\left(\mu - \tfrac{1}{2}\sigma^2\right)t + \sigma W(t)\right].
\]

It follows that $S(t)$ is a \textbf{lognormal process}, since $\log S(t)$ is normally distributed.  
More precisely:
\[
\frac{S(t)}{S(0)} \sim \mathrm{LN}\!\big((\mu - \tfrac{1}{2}\sigma^2)t,\, \sigma^2 t\big),
\]
from which:
\[
\mathbb{E}[S(t)] = S(0)e^{\mu t}, \qquad
\mathrm{Var}[S(t)] = S(0)^2 e^{2\mu t}(e^{\sigma^2 t} - 1).
\]

The process is Markovian and satisfies:
\[
\mathbb{E}[S(t)\mid S(u)] = e^{\mu (t-u)}S(u), \quad t > u.
\]

Along a single trajectory, the average growth rate of the logarithm converges almost surely to:
\[
\frac{1}{t}\log S(t) \to \mu - \frac{1}{2}\sigma^2,
\]
showing that $\mu - \tfrac{1}{2}\sigma^2$ represents the effective \emph{growth rate} observed along a single realization.
\paragraph{Exact Simulation of GBM}

From the closed-form expression above, we can derive an \textbf{exact} procedure for simulating the process at discrete times $0 = t_0 < t_1 < \cdots < t_n$.  
For $S \sim \mathrm{GBM}(\mu,\sigma^2)$, the representation
\[
S(t) = S(0)\exp\!\left[\left(\mu - \tfrac{1}{2}\sigma^2\right)t + \sigma W(t)\right]
\]
implies that, for $u < t$,
\[
S(t) = S(u)\exp\!\left[\left(\mu - \tfrac{1}{2}\sigma^2\right)(t-u) + \sigma(W(t)-W(u))\right].
\]

Since the increments $W(t)-W(u)$ are independent and normally distributed as $\mathcal{N}(0, t-u)$, the following recursive relation holds:
\[
S(t_{i+1}) = S(t_i)\exp\!\left[\left(\mu - \tfrac{1}{2}\sigma^2\right)(t_{i+1}-t_i)
+ \sigma\sqrt{t_{i+1}-t_i}\,Z_{i+1}\right],
\qquad i = 0,1,\dots,n-1,
\]
where $Z_1, Z_2, \dots, Z_n$ are independent standard normal variables.

This method is \textbf{exact} in the sense that the simulated values $(S(t_1), \ldots, S(t_n))$ have exactly the same joint distribution as the continuous-time process $S(t)$ at those discrete times — thus introducing no discretization error.

% Asset/code block omitted pending provenance acceptance.
\paragraph{Time-Dependent Parameters}

In the presence of time-dependent drift and volatility, $\mu(t)$ and $\sigma(t)$, the general expression becomes:
\[
S(t_{i+1}) = S(t_i)\exp\!\left[
\int_{t_i}^{t_{i+1}}\!\!\left(\mu(u) - \tfrac{1}{2}\sigma^2(u)\right)du
+ \int_{t_i}^{t_{i+1}}\!\!\sigma(u)\,dW(u)
\right].
\]
This formulation can be approximated numerically by evaluating the parameters at the midpoint of each interval or by using a deterministic discretization of $\mu(t)$ and $\sigma(t)$.  

In both cases, the method remains consistent with the lognormal structure of the process and ensures strictly positive trajectories, making it one of the most widely used models for simulating financial asset prices.
\paragraph{Risk--Neutral Dynamics}

In pricing applications, the drift \(\mu\) is replaced by the risk--free rate \(r\) under the risk--neutral measure.
If the asset pays a continuous dividend yield \(\delta\), the dynamics become
\[
\frac{dS(t)}{S(t)} = (r - \delta)\,dt + \sigma\,dW(t).
\]
Under this measure, the discounted price \(S(t)e^{-rt}\) is a martingale.
Typical parameterizations include:
\begin{itemize}[leftmargin=10mm]
  \item \textbf{Equity indices:} \(\mu = r - \delta\), with \(\delta\) as the average dividend yield.
  \item \textbf{Exchange rates:} \(\mu = r - r_f\), where \(r_f\) is the foreign interest rate.
  \item \textbf{Commodities:} \(\mu = r - \delta_c\), where \(\delta_c\) may represent storage cost or convenience yield.
\end{itemize}
\paragraph{Path--Dependent Payoffs}

GBM paths are used to evaluate derivative payoffs that depend on the entire price trajectory.  
For discrete monitoring dates \(t_1,\dots,t_n\), an \emph{arithmetic Asian call} has payoff
\[
\Pi = e^{-rT}\max(\bar S - K, 0),
\qquad
\bar S = \frac{1}{n}\sum_{i=1}^{n} S(t_i).
\]
The corresponding Monte Carlo estimator is
\[
\widehat{C}_N = \frac{1}{N}\sum_{j=1}^{N} \Pi^{(j)}.
\]
A \emph{geometric average} option can serve as a variance--reduction control,
since the geometric mean of lognormal prices is itself lognormal and can be priced in closed form.

\begin{remark}[Other path--dependent structures]
The same GBM paths allow simulation of barrier and lookback options by tracking minima, maxima, or barrier crossings along the trajectory.
Continuous monitoring can be approximated by using a sufficiently fine discretization grid.
\end{remark}
\paragraph{Multidimensional Extension}

For \(d\) correlated assets, define \(S(t)=(S_1(t),\dots,S_d(t))^\top\) satisfying
\[
\frac{dS_i(t)}{S_i(t)} = \mu_i\,dt + \sum_{j=1}^{d}A_{ij}\,dW_j(t),
\qquad i=1,\dots,d,
\]
where \(W\) is a \(d\)-dimensional Brownian motion and \(AA^\top=\Sigma\) is the covariance matrix of instantaneous returns.
Discretization on a time grid gives
\[
S_i(t_{k+1}) = S_i(t_k)\,
\exp\!\Big[(\mu_i - \tfrac{1}{2}\sigma_i^2)\Delta t + \sqrt{\Delta t}\,(A Z_k)_i\Big],
\qquad Z_k\sim\mathcal N(0,I_d).
\]

% Asset/code block omitted pending provenance acceptance.



\begin{remark}[Instantaneous covariance]
In this formulation, \(\Sigma\) denotes the \emph{instantaneous covariance} of logarithmic returns:
\(\Sigma = D_\sigma\,\rho\,D_\sigma\), where \(D_\sigma = \mathrm{diag}(\sigma_1,\dots,\sigma_d)\) and \(\rho\) is the correlation matrix.
The subtraction term \(-\tfrac{1}{2}\mathrm{diag}(\Sigma)\Delta t\) corresponds to the Itô correction applied componentwise.
\end{remark}

\begin{remark}[Applications]
GBM provides a consistent basis for pricing and risk analysis of a wide class of derivatives:  
from single--asset options to multi--asset portfolios.  
It also forms the benchmark for more complex models such as stochastic volatility and jump--diffusion processes.
\end{remark}

\paragraph{Asian options (preview).}
For arithmetic Asian options we will use the \emph{geometric Asian} price (available in closed form) as a \emph{control variate}, reusing the same GBM paths (common random numbers) to reduce the variance of the Monte Carlo estimator.
\subsection{Path Simulation: Exact vs Euler--Maruyama (Bias and MSE)} 

Many stochastic differential equations (SDEs) cannot be simulated in closed form.  
In these cases, paths are generated through time discretization of the continuous process.  
The most common approach is the \emph{Euler--Maruyama} method, which provides a simple first–order approximation of the stochastic integral.

\paragraph{Euler scheme.}
Consider the general SDE
\[
dX_t = a(X_t)\,dt + b(X_t)\,dW_t,
\]
where \(a(X_t)\) and \(b(X_t)\) denote drift and diffusion coefficients.  
On a uniform grid \(t_i = i\Delta t\), the discrete update is
\[
\widehat{X}_{i+1} = \widehat{X}_i + a(\widehat{X}_i)\,\Delta t
+ b(\widehat{X}_i)\sqrt{\Delta t}\,Z_i, \qquad Z_i\sim\mathcal N(0,1).
\]
The scheme replaces infinitesimal increments with local linear approximations and becomes exact when \(a\) and \(b\) are constant, as in the geometric Brownian motion.

\paragraph{Convergence and error.}
The accuracy of a discretization is assessed through two main notions:
\begin{itemize}[leftmargin=10mm]
  \item \textbf{Strong convergence:} measures pathwise error  
  \[
  E[|\widehat{X}_T - X_T|^2]^{1/2} = O(\Delta t^{1/2});
  \]
  \item \textbf{Weak convergence:} measures the error on expectations  
  \[
  |E[f(\widehat{X}_T)] - E[f(X_T)]| = O(\Delta t).
  \]
\end{itemize}
The Euler scheme therefore converges strongly with order \(1/2\) and weakly with order \(1\).  
A second–order refinement, known as the \emph{Milstein scheme}, adds the correction term
\[
\tfrac{1}{2}\,b(X_t)b'(X_t)\,(Z^2-1)\,\Delta t,
\]
achieving strong order~1 when \(b\) is smooth.

\paragraph{Bias and mean--square error.}
The total mean--square error (MSE) combines discretization bias and Monte Carlo sampling noise:
\[
\mathrm{MSE} \approx c_1^2\,\Delta t^{2\beta} + \frac{c_2}{n},
\]
where \(\beta\) denotes the weak order of the method and \(n\) the number of paths.  
For the Euler scheme (\(\beta=1\)), halving the global error requires roughly four times more time steps, matching the classical \(O(N^{-1/2})\) convergence of Monte Carlo estimation.

\paragraph{Practical considerations.}
When a closed--form transition is available (for instance in GBM), exact simulation is preferred since it eliminates bias entirely.  
Otherwise, \(\Delta t\) must be chosen small enough that discretization error remains negligible compared to sampling variance.

For path--dependent payoffs, discretization may miss events such as barrier crossings between grid points.  
A \emph{Brownian bridge} interpolation corrects this effect by computing the conditional crossing probability between successive steps:
\[
p = 1 - \exp\!\left[-\frac{2(B - X_i)(B - X_{i+1})}{b_i^2\Delta t}\right],
\]
which can be used to reduce bias in continuously monitored options.

\begin{remark}[Summary]
The Euler--Maruyama scheme provides a reliable and general baseline for the numerical simulation of diffusions.  
It defines the framework for assessing discretization bias and sets the stage for more advanced models such as jump--diffusion processes.
\end{remark}
\subsection{Asian Option Pricing} 

Asian options are the simplest examples of path--dependent payoffs, where the value depends on the entire trajectory of the underlying asset rather than on its terminal level.  
Under the risk--neutral geometric Brownian motion
\[
\frac{dS_t}{S_t} = r\,dt + \sigma\,dW_t,
\]
the asset price evolves as
\[
S(t_{j+1}) = S(t_j)\,
\exp\!\big((r-\tfrac12\sigma^2)\Delta t + \sigma\sqrt{\Delta t}\,Z_{j+1}\big),
\qquad Z_j \sim \mathcal N(0,1).
\]

\paragraph{Payoff definition.}
Let \(0 < t_1 < t_2 < \cdots < t_m = T\) be the observation times and define the discrete arithmetic average
\[
\bar S = \frac{1}{m}\sum_{j=1}^{m} S(t_j).
\]
The discounted payoff of an arithmetic Asian call is
\[
\Pi = e^{-rT}\max(\bar S - K, 0).
\]
For the continuous--time version, the average is replaced by
\[
\bar S_c = \frac{1}{T}\int_0^T S_t\,dt,
\]
which is approached numerically by increasing the number of monitoring dates.

\paragraph{Monte Carlo estimator.}
To estimate the option price, simulate $n$ independent sample paths of \(S_t\) and compute the average payoff:
\[
\widehat{C}_n = \frac{1}{n}\sum_{i=1}^{n}
e^{-rT}\max(\bar S_i - K, 0),
\]
where $\bar S_i$ is the arithmetic mean of prices along path $i$.

% Asset/code block omitted pending provenance acceptance.




\paragraph{Properties.}
The Monte Carlo estimator is unbiased:
\[
\mathbb E[\widehat{C}_n] = C_{\text{true}},
\]
and by the Central Limit Theorem
\[
\sqrt{n}(\widehat{C}_n - C_{\text{true}}) \xrightarrow{d} \mathcal N(0, \sigma_C^2),
\]
where $\sigma_C^2$ is the variance of the discounted payoff.
The standard error thus decreases as \(1/\sqrt{n}\).

\begin{remark}[Geometric vs arithmetic averaging]
The arithmetic average has no closed form.  
However, when the average is geometric,
\[
\bar S_G = \Big(\prod_{j=1}^m S(t_j)\Big)^{1/m},
\]
the payoff \(e^{-rT}\max(\bar S_G - K,0)\) can be evaluated analytically because \(\log \bar S_G\) is normally distributed.  
This provides a useful benchmark to verify the correctness of Monte Carlo implementations.
\end{remark}

\begin{remark}[Discretization]
For sufficiently fine monitoring, the discrete average approximates the continuous one:
\[
\frac{1}{m}\sum_{j=1}^{m} S(t_j) \approx \frac{1}{T}\int_0^T S_t\,dt.
\]
In practice, $m$ is chosen large enough so that discretization bias is negligible relative to sampling error.
\end{remark}

\begin{remark}[Summary]
Asian options illustrate how Monte Carlo methods extend naturally to path--dependent pricing within the geometric Brownian motion framework.  
Their valuation requires only the generation of correlated increments and averaging along each simulated trajectory, 
providing a simple yet powerful setting for studying variance--reduction and quasi--Monte Carlo techniques introduced in the following section.
\end{remark}
\subsection{Extension to Jump--Diffusion Process} 

While pure diffusion models assume continuous paths, real market data often exhibit sudden price movements and heavy tails.  
Jump--diffusion processes extend the classical stochastic framework by superimposing a discontinuous jump component on the continuous diffusion part.

\paragraph{Jump--diffusion process.}
A general jump--diffusion process satisfies
\[
dS_t = \mu S_{t^-}\,dt + \sigma S_{t^-}\,dW_t + S_{t^-}\,dJ_t,
\]
where \(W_t\) is Brownian motion and \(J_t\) is a pure jump process, independent of \(W_t\).  
The notation \(S_{t^-}\) indicates the value immediately before a jump.  
At each jump time \(\tau_j\),
\[
S(\tau_j) = S(\tau_j^-)\,Y_j,
\]
where \(Y_j>0\) is a random multiplicative factor representing the jump magnitude.

If jumps occur according to a Poisson process \(N_t\) with intensity \(\lambda\),
\[
J_t = \sum_{j=1}^{N_t}(Y_j - 1),
\]
then \(S_t\) evolves continuously between jumps and experiences random discontinuities at the jump times.  
The solution of the SDE is
\[
S_t = S_0\,\exp\!\big[(\mu - \tfrac{1}{2}\sigma^2)t + \sigma W_t\big]
\prod_{j=1}^{N_t} Y_j,
\]
which generalizes geometric Brownian motion to a mixture of continuous and jump effects.

\paragraph{Distributional structure.}
Conditional on \(N_t = n\), and if the jumps are lognormal with \(\log Y_j \sim \mathcal N(a,b^2)\),
\[
\log S_t \sim \mathcal N\!\left(\log S_0 + (\mu - \tfrac{1}{2}\sigma^2)t + n\,a,\; \sigma^2 t + n\,b^2\right),
\]
so that \(S_t\) is a Poisson mixture of lognormals.  
This generates leptokurtic distributions and captures the heavy tails typical of financial returns.

\paragraph{Risk--neutral drift adjustment.}
Under the risk--neutral measure, the drift is corrected to ensure that the discounted process \(S_t e^{-(r-\delta)t}\) is a martingale in presence of a continuous dividend yield \(\delta\).  
If \(m = \mathbb E[Y - 1]\) is the mean relative jump size, the compensated drift becomes
\[
\mu = r - \delta - \lambda m.
\]
Then,
\[
\frac{dS_t}{S_{t^-}} = (r-\delta)\,dt + \sigma\,dW_t + (dJ_t - \lambda m\,dt),
\]
so both the diffusion and compensated jump components are martingales.

\paragraph{Simulation on a fixed time grid.}
On a uniform grid \(t_i=i\Delta t\), the logarithm of the price evolves as
\[
X_{i+1} = X_i + (\mu - \tfrac{1}{2}\sigma^2)\Delta t
+ \sigma\sqrt{\Delta t}\,Z_i
+ \sum_{k=1}^{N_i}\log Y_k,
\]
where \(Z_i\sim\mathcal N(0,1)\) and \(N_i\sim\text{Poisson}(\lambda\Delta t)\).  
Exponentiating \(X_i\) yields \(S_i=e^{X_i}\).

% Asset/code block omitted pending provenance acceptance.



This approach aggregates jumps over each time step and provides an efficient discrete--time approximation.  
It is exact when jump arrivals follow a Poisson process independent of the diffusion.

\paragraph{Simulation by jump times.}
Alternatively, the jump times can be generated explicitly.  
Let the interarrival times \(R_j\) follow an exponential distribution with mean \(1/\lambda\).  
Between jumps, the process evolves as a pure geometric Brownian motion.  
When a jump occurs, the price is multiplied by a random factor \(Y_j\).

% Asset/code block omitted pending provenance acceptance.



This method reproduces the exact jump times and is useful for risk metrics sensitive to discontinuities.

\paragraph{Jump--adapted discretization.}
For hybrid SDEs combining diffusions and jumps, an efficient approach is to merge the jump times of the Poisson process with the standard simulation grid.  
Between jumps, the process evolves via the Euler or exact GBM scheme; when a jump occurs, the state is updated as
\[
S(\tau_j) = S(\tau_j^-)\,(1 + \kappa(Y_j)),
\]
where \(\kappa(Y_j)\) defines the jump amplitude.  
This structure preserves the weak order of convergence of the underlying diffusion scheme and is standard in Monte Carlo implementations.

\begin{remark}[Barriers with jumps]
Grid–based simulation may miss crossings caused by instantaneous jumps.  
For barrier payoffs, it is recommended to use a \emph{jump–adapted} grid (the union of regular time steps and jump times) or to simulate jump times explicitly.
\end{remark}


\paragraph{Pure--jump extensions.}
If the Brownian term is removed, the model becomes a pure--jump process, with increments driven entirely by the jump component.  
Examples include the Variance Gamma, Normal Inverse Gaussian, and other Lévy processes, which can be viewed as random time--changes of Brownian motion.  
These models exhibit infinite activity and offer flexible shapes for return distributions.

\begin{remark}[Applications]
Jump--diffusion models capture abrupt price movements and heavy tails observed in financial data.  
They are used for option pricing, stress testing, and portfolio risk simulations, bridging the gap between continuous Gaussian models and empirical market behavior.
\end{remark}

\paragraph{Consistency checks.}
On synthetic paths we verify: (i) for BM, \( \hat{\mathbb E}[W(t)]\approx0\) and \(\widehat{\mathrm{Var}}[W(t)]\approx t\);
(ii) for GBM, \( \hat{\mathbb E}[S_t]\approx S_0 e^{\mu t}\);
(iii) in \(d\) dimensions, the empirical correlation matrix of log-returns matches the target \(\Sigma\) within \(O(N^{-1/2})\) fluctuations.

% ---- next approved source slice ----

\section{Variance Reduction Techniques} 

The previous sections introduced the foundations of Monte Carlo simulation, random number generation, and stochastic path modeling for asset prices. While these methods provide unbiased estimators for expectations and option prices, their statistical accuracy improves only at the slow rate \(O(N^{-1/2})\). In practical financial applications—such as pricing complex derivatives or computing portfolio risk metrics—this convergence can be prohibitively expensive in terms of computational cost. Variance reduction techniques address this limitation by increasing estimator efficiency: they aim to achieve the same level of accuracy with fewer simulations, or equivalently, higher precision for the same computational effort. These methods modify either the structure of the estimator or the generation of random samples, while preserving unbiasedness and consistency.

This section presents the main variance reduction strategies used in quantitative finance. The first group, based on analytical exploitation, uses known relationships or auxiliary quantities with known expectations to stabilize results—most notably the method of control variates. The second group, based on sampling modification, alters the input randomness to reduce variability without affecting the mean—this includes antithetic variates, stratified sampling, and quasi–Monte Carlo sequences. Each technique is introduced from first principles, with mathematical formulation, efficiency analysis, and numerical implementation. Applications focus on European and Asian options, comparing plain Monte Carlo estimators with their variance–reduced and quasi–Monte Carlo counterparts.

Independent random draws are not always optimal: introducing controlled dependence or structure across samples can significantly reduce overall estimator variance without altering its expectation. Although this breaks the assumption of independence among simulated observations, the resulting estimators remain unbiased, and confidence intervals can still be derived using the same asymptotic principles, provided that variance estimation accounts for the induced correlation. Variance reduction therefore represents a natural and essential evolution of the Monte Carlo method, linking theoretical convergence with practical computational efficiency and forming the basis for advanced techniques such as quasi–Monte Carlo methods discussed later in this section.
\subsection{Control Variates} 

The method of control variates is one of the most effective and widely applicable techniques for improving the efficiency of Monte Carlo simulation. It exploits the information contained in quantities with known expectations to reduce the variance of an estimator for an unknown expectation.  

\begin{definition}[Control variate estimator] 
Let \((X_i, Y_i)\), \(i = 1, \dots, n\), be i.i.d. pairs of random variables, where \(Y_i\) represents the output of interest and \(X_i\) is an auxiliary variable (the \emph{control}) whose expectation \(\mathbb{E}[X]\) is known.  
The standard Monte Carlo estimator of \(\mathbb{E}[Y]\) is the sample mean \(\bar{Y}\).  
The control variate estimator introduces a correction term based on the observed deviation of \(X_i\) from its known mean:
\[
\widehat{Y}(b) = \bar{Y} - b(\bar{X} - \mathbb{E}[X]).
\]
For any fixed coefficient \(b\), the estimator remains unbiased since \(\mathbb{E}[\widehat{Y}(b)] = \mathbb{E}[Y]\).
\end{definition}

\begin{remark}[Variance and optimal coefficient]
The variance per replication is
\[
\mathrm{Var}[Y - b(X - \mathbb{E}[X])] = \sigma_Y^2 - 2b\sigma_X\sigma_Y\rho_{XY} + b^2\sigma_X^2,
\]
where \(\sigma_X^2 = \mathrm{Var}[X]\), \(\sigma_Y^2 = \mathrm{Var}[Y]\), and \(\rho_{XY}\) is the correlation between \(X\) and \(Y\).  
The variance is minimized for
\[
b^\ast = \frac{\mathrm{Cov}(X,Y)}{\mathrm{Var}(X)},
\]
yielding
\[
\mathrm{Var}[\widehat{Y}(b^\ast)] = (1 - \rho_{XY}^2)\,\mathrm{Var}[\bar{Y}],
\]
which shows that the effectiveness of the method depends entirely on the strength of correlation between the control and the quantity of interest.
\end{remark}

\begin{remark}[Estimation of the coefficient]
In practice, the optimal coefficient is estimated from simulated data as
\[
\hat{b}_n = 
\frac{\sum_{i=1}^n (X_i - \bar{X})(Y_i - \bar{Y})}
     {\sum_{i=1}^n (X_i - \bar{X})^2},
\]
which converges almost surely to \(b^\ast\) as \(n \to \infty\).  
Using the estimated coefficient,
\[
\widehat{Y}(\hat{b}_n) = \bar{Y} - \hat{b}_n(\bar{X} - \mathbb{E}[X]),
\]
produces an estimator that is asymptotically unbiased and achieves the same variance reduction as when using \(b^\ast\).
\end{remark}

\begin{remark}[Multiple control variates]
The method generalizes naturally to multiple controls.  
If \(X = (X^{(1)},\dots,X^{(d)})^\top\) is a vector of control variables with known expectations, the optimal coefficient vector is
\[
b^\ast = \Sigma_X^{-1}\Sigma_{XY},
\]
where \(\Sigma_X = \mathrm{Cov}(X)\) and \(\Sigma_{XY} = \mathrm{Cov}(X,Y)\).  
The corresponding variance reduction factor is \(1 - R^2\), with
\[
R^2 = \frac{\Sigma_{XY}^\top \Sigma_X^{-1}\Sigma_{XY}}{\sigma_Y^2},
\]
representing the fraction of total variance eliminated by the control variates.
\end{remark}

\begin{remark}[Interpretation]
The control variate method can be viewed as a linear regression adjustment in which the residuals \(Y - b^\ast(X - \mathbb{E}[X])\) are uncorrelated with the controls.  
The closer the payoff \(Y\) can be expressed as a linear combination of the control variables, the greater the resulting variance reduction.  
When the correlation \(|\rho_{XY}|\) approaches one, the reduction in variance can exceed an order of magnitude compared to plain Monte Carlo estimation.
\end{remark}
\paragraph{Single Control Variate: Python Implementation}

The following pseudocode illustrates a practical Python implementation of the single control variate estimator described in Definition~\sourceref{def:control-variate}.
Given simulation outputs \(Y_i\) and a control variate \(X_i\) with known expectation \(\mathbb{E}[X]\), the estimator computes the optimal coefficient 
\(\hat{b}\) and produces the variance–reduced estimate
\[
\hat{Y}_{cv} = \frac{1}{n}\sum_{i=1}^n \left( Y_i - \hat{b}\,\left( X_i - \mathbb{E}[X] \right) \right).
\]

% Asset/code block omitted pending provenance acceptance.
\subsection{Antithetic Variates} 

The method of antithetic variates reduces variance by introducing negative dependence between paired replications of a simulation. Instead of generating independent samples, each random input is combined with its \emph{antithetic} counterpart, constructed to offset the stochastic fluctuations of the original sample. The approach relies on the symmetry of random number transformations and the monotonicity of the simulation function with respect to its inputs.

\begin{definition}[Antithetic pair and estimator]
Let \(U \sim \mathcal{U}(0,1)\). Since \(1-U\) has the same distribution, the pair \((U,1-U)\) forms an \emph{antithetic pair}.  
Similarly, if \(Z \sim \mathcal{N}(0,1)\), then \((-Z)\) is an antithetic normal variate.  
If the simulation output is \(Y = f(U)\) or \(Y = f(Z)\), the corresponding antithetic output is \(\widetilde{Y} = f(1-U)\) or \(f(-Z)\).  
The antithetic variates estimator of \(\mathbb{E}[Y]\) is defined as
\[
\widehat{Y}_{AV} = \frac{1}{n}\sum_{i=1}^{n} \frac{Y_i + \widetilde{Y}_i}{2},
\]
where the pairs \((Y_i,\widetilde{Y}_i)\) are i.i.d.\ and each pair shares the same marginal distribution.
\end{definition}

\begin{remark}[Unbiasedness and variance]
The estimator \(\widehat{Y}_{AV}\) is unbiased since both \(Y_i\) and \(\widetilde{Y}_i\) have expectation \(\mathbb{E}[Y]\).  
Its asymptotic distribution follows from the Central Limit Theorem:
\[
\sqrt{n}\,\frac{\widehat{Y}_{AV}-\mathbb{E}[Y]}{\sigma_{AV}} \xrightarrow{d} \mathcal{N}(0,1),
\]
where
\[
\sigma_{AV}^2 = \mathrm{Var}\!\left[\frac{Y + \widetilde{Y}}{2}\right]
= \tfrac{1}{2}\mathrm{Var}[Y] + \tfrac{1}{2}\mathrm{Cov}[Y,\widetilde{Y}].
\]
Variance reduction occurs whenever \(\mathrm{Cov}[Y,\widetilde{Y}] < 0\). The larger the negative covariance between the paired outputs, the greater the reduction in sampling variance.
\end{remark}

\begin{remark}[Condition for variance reduction]
Comparing with a standard Monte Carlo estimator based on \(2n\) independent samples, antithetic sampling is more efficient if
\[
\mathrm{Var}[Y_i + \widetilde{Y}_i] < 2\,\mathrm{Var}[Y_i],
\]
which simplifies to
\[
\mathrm{Cov}[Y_i,\widetilde{Y}_i] < 0.
\]
This condition requires that the simulation mapping \(f(\cdot)\) from inputs to outputs transforms positively correlated antithetic inputs into negatively correlated outputs. In particular, if \(Y = f(Z)\) is monotone increasing in each input \(Z_j\), and the antithetic counterpart is defined by \(\widetilde{Y} = f(-Z)\), then the covariance is nonpositive, ensuring a variance reduction.
\end{remark}

\begin{remark}[Variance decomposition]
When the simulation depends on a symmetric input distribution, such as \(Z \sim \mathcal{N}(0,I)\), the integrand can be decomposed into symmetric and antisymmetric parts:
\[
f_0(z) = \tfrac{1}{2}[f(z) + f(-z)], \qquad
f_1(z) = \tfrac{1}{2}[f(z) - f(-z)].
\]
These satisfy \(\mathbb{E}[f_0(Z)f_1(Z)] = 0\) and
\[
\mathrm{Var}[f(Z)] = \mathrm{Var}[f_0(Z)] + \mathrm{Var}[f_1(Z)].
\]
Antithetic sampling eliminates the contribution of the antisymmetric component \(f_1(Z)\), which represents the variance due to the asymmetry of the integrand. It is therefore most effective when \(f\) is nearly linear or weakly nonlinear in its inputs.
\end{remark}

\begin{remark}[Systematic extensions]
The antithetic transformation \(Z \mapsto -Z\) can be generalized to any orthogonal transformation \(T\) satisfying \(TZ \sim \mathcal{N}(0,I)\).  
Applying multiple transformations \(T^1Z,\dots,T^kZ\) defines the systematic estimator
\[
f_0(Z) = \frac{1}{k}\sum_{i=1}^{k} f(T^iZ),
\]
which preserves the expected value and achieves variance reduction when
\[
\sum_{j=1}^{k-1}\mathrm{Cov}[f(Z),f(T^jZ)] < 0.
\]
The classical antithetic case corresponds to \(k=2\) and \(T=-I\).
\end{remark}

\begin{remark}[Practical advantages]
The antithetic variates method is easy to implement and computationally inexpensive, as it reuses the same random inputs with opposite signs or complementary values. It preserves unbiasedness, introduces only minimal correlation structure, and is particularly effective for approximately monotone or symmetric payoffs. The method provides variance reduction without altering the sampling distribution or requiring auxiliary expectations.
\end{remark}

% ---- next approved source slice ----

\subsection{Quasi--Monte Carlo Methods}

Quasi--Monte Carlo (QMC) methods represent one of the most significant developments in modern numerical simulation. They combine the conceptual structure of Monte Carlo integration with a deterministic construction of sampling points, achieving a substantially faster convergence rate than classical pseudo--random techniques. The core idea is to replace independent pseudo--random numbers with \emph{low--discrepancy sequences}, designed to cover the hypercube $[0,1]^d$ with exceptional uniformity.

The motivation behind this approach is that the error rate of standard Monte Carlo is limited by the statistical convergence $O(N^{-1/2})$, which follows from the Central Limit Theorem. In contrast, for low--discrepancy sequences, one can obtain nearly linear convergence rates, typically of the form
\[
O\!\left( \frac{(\log N)^d}{N} \right),
\]
which, for moderate values of $d$, results in a drastic improvement in precision for the same computational cost.  
In path--dependent and multi--asset financial problems, this acceleration has considerable practical importance.
\paragraph{Discrepancy and point uniformity}

The notion of \emph{discrepancy} is central to understanding the effectiveness of QMC.  
Given a sequence of points $\{x_1,\dots,x_N\}\subset [0,1]^d$, the \emph{star discrepancy} $D_N^*$ is defined as
\[
D_N^* := 
\sup_{u \in [0,1]^d}
\left|
\frac{1}{N}
\sum_{i=1}^N 1_{\{x_i\in [0,u]\}}
- \prod_{j=1}^d u_j
\right|.
\]
It measures the maximum deviation between the empirical distribution of points in axis-aligned boxes anchored at the origin and their theoretical volumes.  
Low discrepancy implies that the points \emph{fill} the space very uniformly, avoiding clusters and empty regions typical of pseudo--random sequences.

For independent pseudo--random points, one typically has
\[
\mathbb{E}[D_N^*] = O(N^{-1/2}),
\]
whereas there exist deterministic sequences satisfying
\[
D_N^* = O\!\left(\frac{(\log N)^d}{N}\right),
\]
which is essentially optimal for high-dimensional numerical integration.
\paragraph{The Koksma--Hlawka inequality}

The integration error of QMC estimators is governed by the Koksma--Hlawka inequality.  
If $f:[0,1]^d\to\mathbb{R}$ has bounded variation in the sense of Hardy--Krause, then
\[
\big|I(f) - \widehat{I}_N\big|
\le V_{\mathrm{HK}}(f)\, D_N^*,
\qquad 
I(f)=\int_{[0,1]^d} f(u) \, du,
\qquad 
\widehat{I}_N=\frac{1}{N}\sum_{i=1}^N f(x_i),
\]
where $V_{\mathrm{HK}}(f)$ measures the structural complexity of $f$.

This result is fundamental for two reasons:
\begin{enumerate}
    \item it provides a \emph{deterministic} error bound, unlike classical Monte Carlo which offers only probabilistic guarantees;
    \item it shows that the convergence rate is driven entirely by the discrepancy of the chosen sequence.
\end{enumerate}
Although the bound is often pessimistic and the variation $V_{\mathrm{HK}}(f)$ is rarely computable in practice, the inequality provides the theoretical justification for QMC.
\paragraph{Low--discrepancy sequences}

The choice of sequence is crucial in QMC. The most common ones include:

\paragraph{Halton sequence.}
Constructed using prime bases, it generalizes the one-dimensional van der Corput sequence.  
It is simple to implement but suffers from strong correlations across dimensions when $d>10$--$20$, making it unsuitable for most financial applications.

\paragraph{Sobol' sequence.}
The standard in quantitative finance.  
It is built from primitive polynomials over $\mathbb{F}_2$ and direction matrices, using bitwise XOR operations to generate each point.  
Its key properties include:
\begin{itemize}
    \item excellent uniformity in 2D projections;
    \item stable behaviour even for dimensions in the hundreds or thousands;
    \item natural compatibility with dimension reduction techniques such as Brownian Bridge and PCA.
\end{itemize}
The $i$-th point is generated via a Gray code expansion and linear operations over $\mathbb{F}_2$, ensuring that small changes in the index produce nearby points, improving global space-filling behaviour.

\paragraph{Faure and Niederreiter--Xing sequences.}
Based on more sophisticated constructions (finite fields, algebraic curves), they offer extremely low discrepancy in theory, but are less commonly used in finance due to implementation complexity.
\paragraph{Error estimation via randomization}

Because QMC is deterministic, it does not inherently provide a statistical error estimator.  
This limitation is addressed through \emph{randomized QMC} (RQMC).  
Common randomization techniques include:
\begin{itemize}
    \item \textbf{Digital shift}: bitwise addition of a random number in $\mathbb{F}_2$;
    \item \textbf{Owen scrambling}: hierarchical random permutations of bits while preserving low discrepancy in expectation;
    \item \textbf{Cranley–Patterson rotation}: random toroidal shifts.
\end{itemize}

Randomization generates $K$ independent versions of the low--discrepancy sequence:
\[
\widehat{I}_N^{(k)} = \frac{1}{N}\sum_{i=1}^N f\bigl(x_i^{(k)}\bigr),
\qquad k=1,\dots,K,
\]
and the overall estimator
\[
\overline{I}_N = \frac{1}{K}\sum_{k=1}^K \widehat{I}_N^{(k)}
\]
is unbiased.  
The empirical variance across the $K$ samples provides a statistical error estimate, typically with variance
\[
\operatorname{Var}(\overline{I}_N) = O\!\left(\frac{1}{N^2}\right),
\]
in contrast with the $O(1/N)$ of classical Monte Carlo.
\paragraph{Convergence and effective dimension}

In high dimensions, theoretical bounds may deteriorate due to the \emph{curse of dimensionality}.  
However, many financial problems have exploitable structure.  
In particular, Brownian motion allows a substantial reduction of the \emph{effective dimension} through:
\begin{itemize}
    \item \textbf{Brownian Bridge construction}: concentrates the largest variances in early coordinates of the QMC sequence;
    \item \textbf{PCA}: diagonalizes the covariance matrix of log-returns, reducing complexity to a few dominant factors.
\end{itemize}
Both techniques substantially enhance QMC performance, preserving its near-linear convergence rate.
\paragraph{Financial applications and numerical implementation}

Quasi--Monte Carlo methods have become a standard tool for the numerical valuation of complex derivatives and the simulation of high--dimensional stochastic models in quantitative finance. Their main advantage lies in the substantial variance reduction and the increased stability of numerical estimates, especially in problems where classical Monte Carlo suffers from slow convergence or high noise.

Typical application domains include:
\begin{itemize}
    \item \emph{Path--dependent options}, such as arithmetic Asian options, barrier options, lookback options and cliquet structures, where the payoff depends on the whole trajectory of the underlying asset;
    \item \emph{Multi--asset derivatives}, e.g.\ basket options, quanto structures and correlation products, where one has to simulate a vector of correlated risk factors over many time steps;
    \item \emph{Risk measurement}, such as the computation of Value--at--Risk (VaR) and Expected Shortfall (ES) for large portfolios under complex scenario generation schemes.
\end{itemize}

In all these settings, the underlying numerical task can be formulated as the approximation of a high--dimensional expectation
\[
\mathbb{E}[f(X)] = \int_{\mathbb{R}^d} f(x)\,\psi(x)\,dx,
\]
where $X$ encodes the relevant risk factors (e.g.\ discretised Brownian paths or multi--asset returns) and $\psi$ denotes their joint distribution. QMC methods tackle this problem by mapping low--discrepancy points $u_i \in [0,1]^d$ into samples $x_i$ via suitable transformations (inverse CDFs, copula constructions, Brownian Bridge, PCA, etc.), and then forming the estimator
\[
\widehat{I}_N^{\mathrm{QMC}} = \frac{1}{N} \sum_{i=1}^N f(x_i).
\]

Compared to standard Monte Carlo, practitioners typically observe:
\begin{itemize}
    \item a much smoother convergence of prices and Greeks as a function of $N$;
    \item a significant reduction in the number of paths required to reach a given error tolerance;
    \item more stable hedging parameters, which is crucial for risk management and P\&L attribution.
\end{itemize}

\paragraph{Example: QMC pricing of an arithmetic Asian call}

We illustrate a concrete implementation of QMC in the pricing of an arithmetic Asian call option under the risk--neutral Black--Scholes model. Let $(S_t)_{t\in[0,T]}$ follow the geometric Brownian motion
\[
dS_t = r S_t\,dt + \sigma S_t\,dW_t,
\]
with initial spot $S_0$, risk--free rate $r$ and volatility $\sigma$.  
Consider an arithmetic Asian call with strike $K$ and maturity $T$, monitored at $m$ equally spaced times $0 < t_1 < \dots < t_m = T$. Its discounted payoff is
\[
\Pi = e^{-rT}\,\max\!\left( \frac{1}{m}\sum_{j=1}^m S_{t_j} - K,\;0 \right).
\]

In a QMC framework, we proceed as follows:
\begin{enumerate}
    \item generate $N$ points $u_i \in [0,1]^m$ using a Sobol' low--discrepancy sequence;
    \item transform each component via the inverse standard normal CDF to obtain $Z_i \in \mathbb{R}^m$ with i.i.d.\ $N(0,1)$ entries;
    \item construct GBM paths using the exact discretisation
    \[
    S_{t_{j+1}} = S_{t_j} \exp\!\left( \bigl(r - \tfrac{1}{2}\sigma^2\bigr)\Delta t + \sigma \sqrt{\Delta t}\,Z_{i,j} \right), 
    \quad \Delta t = \frac{T}{m};
    \]
    \item compute the pathwise payoff $\Pi_i$ and the QMC estimator
    \[
    \widehat{C}^{\mathrm{QMC}} = \frac{1}{N}\sum_{i=1}^N \Pi_i.
    \]
\end{enumerate}

The Python function below shows a concrete implementation of this procedure using the Sobol' engine from \texttt{scipy.stats.qmc}. Although QMC is deterministic, we also report a standard--error estimate obtained from the sample variance of the payoffs, which is useful in practice as a heuristic error indicator.

% Asset/code block omitted pending provenance acceptance.



In a production--grade implementation, this QMC engine can be embedded into a larger pricing framework, combined with Brownian Bridge or PCA to reduce the effective dimension and extended to multi--asset settings by introducing correlated normal variates through Cholesky factorisation of the covariance matrix.  
Such architectures leverage the superior convergence of QMC to obtain more accurate and stable prices and Greeks, often with one or two orders of magnitude fewer paths than standard Monte Carlo.
\subsubsection{Comparative Numerical Benchmark for Variance Reduction Methods}


To compare the practical effectiveness of the variance reduction methods introduced in this chapter, we consider a simple but representative benchmark problem: the pricing of a European call option under the Black--Scholes model. The goal is not to study model complexity, but rather to isolate the effect of each simulation technique on the statistical accuracy of the Monte Carlo estimator.

Under the risk--neutral measure, the terminal stock price satisfies
\[
S_T = S_0 \exp\!\left[\left(r - \tfrac{1}{2}\sigma^2\right)T + \sigma \sqrt{T}\, Z\right],
\qquad Z \sim \mathcal{N}(0,1),
\]
and the option price is
\[
C_0 = e^{-rT}\,\mathbb{E}\!\left[(S_T-K)^+\right].
\]
The corresponding crude Monte Carlo estimator is
\[
\widehat{C}_N
=
e^{-rT}\,\frac{1}{N}\sum_{i=1}^{N}(S_T^{(i)}-K)^+.
\]

In the numerical experiment below, this same pricing problem is evaluated using several estimators: crude Monte Carlo, antithetic variates, control variates, stratified sampling, importance sampling, quasi--Monte Carlo, and selected combinations of these techniques. Since the Black--Scholes price is known in closed form, it provides a natural benchmark against which the numerical estimates can be compared.

The first comparison concerns the convergence of the estimators as the number of simulations increases, together with the decay of the absolute pricing error relative to the analytical Black--Scholes value.

% Asset/code block omitted pending provenance acceptance.


Figure~\sourceref{fig:vr_benchmark_convergence} shows that all estimators converge toward the same theoretical option price, but not at the same practical speed. Variance reduction does not change the expectation being estimated; rather, it reduces the dispersion of the estimator around the true value. As a consequence, some methods achieve a visibly smaller error for the same computational budget.

A second comparison focuses on the variability of the estimators themselves. This is measured through the empirical standard deviation across repeated runs and through the implied variance reduction relative to crude Monte Carlo.

% Asset/code block omitted pending provenance acceptance.


Figure~\sourceref{fig:vr_benchmark_variability} highlights the central quantitative point of the chapter: the efficiency gain of a Monte Carlo method is measured by how much it reduces estimator variance while preserving unbiasedness or asymptotic correctness. In this benchmark, control variates and quasi--Monte Carlo methods typically provide the strongest improvement, while antithetic and stratified sampling also yield meaningful gains relative to plain simulation.



\begin{remark}
This benchmark should be interpreted as an illustration of relative efficiency rather than as a universal ranking. The performance of a variance reduction technique depends on the structure of the payoff, the dimension of the problem, and the degree of smoothness of the integrand. Nevertheless, even in this simple Black--Scholes setting, the comparison makes clear that variance reduction can produce substantial gains over crude Monte Carlo, thereby justifying its central role in practical computational finance.
\end{remark}
